\documentclass[11pt]{article}

\usepackage{fontspec}
\setmainfont{cmunrm.otf}[BoldFont=cmunbx.otf,ItalicFont=cmunti.otf,BoldItalicFont=cmunbi.otf]
\usepackage[english,main=russian]{babel}

\usepackage[margin=1in]{geometry}
\usepackage{amsmath,amssymb,amsthm,mathtools}
\usepackage{booktabs}
\usepackage{enumitem}
\usepackage{xcolor}
\usepackage[colorlinks=true,linkcolor=black,citecolor=black,urlcolor=blue]{hyperref}

\theoremstyle{plain}
\newtheorem{theorem}{Теорема}
\newtheorem{proposition}{Предложение}
\theoremstyle{definition}
\newtheorem{definition}{Определение}
\newtheorem{axiom}{Аксиома}
\theoremstyle{remark}
\newtheorem{remark}{Замечание}

\newcommand{\V}{\mathrm{V}}
\newcommand{\R}{\mathrm{R}}
\newcommand{\confidence}{\operatorname{confidence}}
\newcommand{\Supply}{\operatorname{Supply}}
\newcommand{\VerifiedValue}{\operatorname{VerifiedValue}}
\newcommand{\ExtRev}{\mathrm{ExtRev}}
\newcommand{\Turnover}{\operatorname{Turnover}}

\title{\textbf{Kredo}: экономика взаимного кредита на основе вклада\\
с непередаваемой (soulbound) репутацией и компенсаторной эмиссией}

\author{Fedor Chishchin\\
\small Независимый исследователь\\
\small \texttt{fedorstartup@gmail.com}}

\date{\today}

\begin{document}
\maketitle

\begin{abstract}
Мы представляем \emph{Kredo} --- экономический механизм, спроектированный так, чтобы вознаграждать создание ценности, а не позицию владения. Устройство выводится сверху вниз: пять философских аксиом формализуются в семь инвариантов состояния, и каждая функция перехода обязана их сохранять. В экономике используются два актива: передаваемый токен вклада $\V$ (рабочие деньги и право на долю внешней прибыли клуба) и непередаваемый «soulbound» токен репутации $\R$. Кредит выпускается в момент подтверждённой транзакции и сопровождается \emph{компенсаторной эмиссией} --- отложенным, помещённым в эскроу «дивидендом всем», величина которого управляется динамическим коэффициентом $\varepsilon$, убывающим по мере роста наблюдаемой ставки дефолтов. Мы доказываем, что этот коэффициент удерживает денежную массу ограниченной относительно подтверждённой ценности при ненулевой ставке дефолтов (Теорема~\ref{thm:stability}) и что операции с фондом ликвидности нейтральны к курсу (Теорема~\ref{thm:neutrality}), так что цена токена реагирует только на фундаментальные величины. Управление использует квадратичное ($\sqrt{\R}$) голосование при двойном кворуме репутации и доли, что вместе с непередаваемостью $\R$ делает захват капиталом или Sybil-аккаунтами экономически невыгодным. Мы выводим минимальное условие выживания по внешней выручке для неубывающего курса. Детерминированный агентный симулятор, реализующий полный механизм, прогнан на $315$ запусках устойчивости --- Монте-Карло, развёртка из $192$ точек по параметрам и стресс-тесты с комбинированными атаками --- ни один из которых не привёл к экономическому коллапсу, а отдельный анализ режимов воспроизводит предсказанные траектории курса. Весь код, выводы и результаты выложены в открытый доступ.
\end{abstract}

\noindent\textbf{Ключевые слова:} взаимный кредит; комплементарная валюта; проектирование механизмов; токеномика; квадратичное голосование; устойчивость к Sybil-атаке; агентное моделирование.

\section{Введение}

Два широко распространённых класса экономической организации несут в себе структурную несправедливость. В \emph{рентном капитализме} держатель капитала зарабатывает, не внося вклада: ценность извлекается из позиции владения, а не создаётся, и создатель ценности и получатель прибыли часто оказываются разными агентами. В \emph{эгалитарных кооперативах и коммунах} все участники равны, но отсутствие механизма роста и вознаграждения порождает классическую проблему безбилетника: вклад не вознаграждается дифференцированно, активные участники несут непропорциональные издержки, а система стагнирует.

Эта работа изучает механизм \emph{Kredo}, задуманный как третий вариант: вознаграждение пропорционально созданию ценности, при котором активные участники зарабатывают больше пассивных, но никто не остаётся ни с чем, и при котором система защищается от инфляции, паники и мошенничества без дискреционного вмешательства.

Наш вклад состоит в следующем.
\begin{enumerate}[leftmargin=*,itemsep=2pt]
    \item Метод проектирования сверху вниз: пять аксиом (раздел~\ref{sec:axioms}) формализуются в семь инвариантов, которые обязана сохранять каждая функция перехода (раздел~\ref{sec:invariants}), что даёт спецификацию, в которой заложенная этика становится машинно-проверяемым свойством.
    \item Механизм компенсаторной эмиссии с динамическим коэффициентом $\varepsilon$ и доказательство того, что он ограничивает массу относительно подтверждённой ценности при ненулевых дефолтах (Теорема~\ref{thm:stability}).
    \item Правило ценообразования, строго привязанное к фундаментальным величинам, с доказательством нейтральности операций фонда к курсу (Теорема~\ref{thm:neutrality}), а также механизм дисконтированного вывода для устойчивости к банк-рану.
    \item Управленческая конструкция, сочетающая $\sqrt{\R}$-голосование с двойным кворумом и непередаваемым токеном репутации, и аргумент о невыгодности захвата как капиталом, так и Sybil-аккаунтами (Предложение~\ref{prop:sybil}).
    \item Минимальное условие выживания по внешней выручке (Предложение~\ref{prop:survival}) и эмпирическая валидация на $315$ симуляционных запусках.
\end{enumerate}

Все артефакты --- механизм, его выводы, симулятор, протоколы экспериментов и сырые результаты --- находятся в открытом доступе.\footnote{\url{https://github.com/fedorello/kredo-model}}

\section{Смежные работы}\label{sec:related}

Идея о том, что деньги не должны приносить ренту праздным держателям, восходит к предложению Гезеля о плате за хранение валюты \cite{gesell1916}. Системы взаимного кредита, в которых денежная масса создаётся эндогенно двусторонней торговлей и в сумме равна нулю, имеют долгую практическую историю \cite{greco2009}; эмпирические работы о швейцарском WIR/\emph{Wirtschaftsring} документируют их контрциклическое, макростабилизирующее поведение \cite{stodder2009}. Современные блокчейн-комплементарные валюты, такие как Circles~UBI \cite{circles2020} и GoodDollar \cite{gooddollar2021}, распределяют вновь выпущенные единицы широко; «дивиденд всем» в Kredo --- это отложенный, взвешенный по вкладу и обусловленный дефолтами вариант этой идеи. Бездоверительный выпуск и отказ от premine следуют дизайн-этосу пермиссионлесс-систем ценности \cite{nakamoto2008}. Наше управление опирается на квадратичное голосование \cite{lalley2018} и квадратичное финансирование \cite{buterin2019}, где взвешивание по квадратному корню ограничивает плутократическое доминирование; именно непередаваемый (soulbound) характер репутации делает квадратичное взвешивание устойчивым к покупке голосов. Устойчивость к Sybil-атакам анализируется в традиции Дусёра \cite{douceur2002}. Наконец, мы оцениваем механизм агентным моделированием \cite{bonabeau2002} --- стандартным инструментом изучения эмерджентного поведения множества взаимодействующих агентов там, где аналитический разбор в замкнутой форме невозможен.

\section{Модель}

\subsection{Аксиомы}\label{sec:axioms}

\begin{axiom}[Ценность создаётся, а не извлекается]\label{ax:create}
Вознаграждение должно соответствовать созданию ценности, а не позиции владения.
\end{axiom}
\begin{axiom}[Ценность субъективна, но согласованна]\label{ax:price}
Справедливая цена --- та, о которой две стороны договариваются добровольно при наличии альтернатив.
\end{axiom}
\begin{axiom}[Симметрия выигрыша и проигрыша]\label{ax:symmetry}
Все участники структурно связаны судьбой системы; ни один класс не защищён от рисков остальных.
\end{axiom}
\begin{axiom}[Отсутствие привилегированного знания]\label{ax:transparency}
Все правила, балансы и решения прозрачны.
\end{axiom}
\begin{axiom}[Власть от вклада]\label{ax:power}
Влияние на правила вытекает из доказанного полезного участия, а не из накопленного капитала.
\end{axiom}

\subsection{Состояние}

\begin{definition}[Состояние системы]
В дискретный момент $t$ состояние --- это кортеж
\[
S_t = (M_t,\ b_t,\ r_t,\ d_t,\ F_t,\ E_t,\ G_t,\ H_t),
\]
где $M_t$ --- конечное множество участников; $b_t: M_t \to \mathbb{R}$ --- баланс $\V$ (может быть отрицательным вплоть до кредитного лимита); $r_t: M_t \to \mathbb{R}_{\ge 0}$ --- репутация; $d_t: M_t \to \mathbb{R}_{\ge 0}$ --- накопленный долг; $F_t \in \mathbb{R}_{\ge 0}$ --- фонд ликвидности (во внешнем стабильном активе, USDC); $E_t \in \mathbb{R}_{\ge 0}$ --- эскроу распределения; $G_t \in \mathbb{R}_{\ge 0}$ --- генезис-пул; $H_t$ --- неизменяемая (append-only) история транзакций.
\end{definition}

\subsection{Инварианты}\label{sec:invariants}

Семь инвариантов ниже формализуют Аксиомы~\ref{ax:create}--\ref{ax:power}. Пусть $N_{\mathrm{credit}}(\tau)$ --- вновь выпущенная (кредитная) часть транзакции $\tau$, $\confidence(\tau)\in[0,1]$ --- оценка доверия из раздела~\ref{sec:confidence}, а $C(\tau) = N_{\mathrm{credit}}(\tau)\,\confidence(\tau)$ --- её \emph{подтверждённый вклад}. Обозначим
\[
\Supply(t) = \sum_{m} \max\bigl(0,b_t(m)\bigr) + E_t + G_t, \qquad
\VerifiedValue(t) = \sum_{\tau \in H_t}\bigl(1+\varepsilon(\tau)\bigr)C(\tau) \;+\; \Gamma_t \;-\; \mathrm{Burned}(t),
\]
где $\varepsilon(\tau)$ --- коэффициент компенсации из \eqref{eq:epsilon}, зафиксированный в момент $\tau$, $\Gamma_t$ --- накопленные генезис-гранты (сами финансируемые из реализованной внешней выручки и штрафов, раздел~\ref{sec:emission}), а $\mathrm{Burned}(t)$ --- суммарный объём, удалённый дефолтами, штрафами и конвертациями. Поскольку $\Supply$ содержит выпущенный принципал, эскроу-компенсацию и генезис-гранты за вычетом сожжённого, $\VerifiedValue$ определён так, чтобы совпадать с ним; любой остаточный зазор в (I3) отражает лишь задержку между мошенничеством и его обнаружением.

\begin{description}[leftmargin=*,itemsep=3pt]
    \item[(I1) Сохранение при переводах.] Для любого чистого перевода $\sum_{m} b_{t}(m) = \sum_{m} b_{t+1}(m)$.
    \item[(I2) Эмиссия под подтверждение.] $\V$ выпускается только внутри протокола подтверждения и только если $\confidence(\tau) \ge \theta_{\min}$.
    \item[(I3) Масса отслеживает ценность.] $\bigl|\Supply(t) - \VerifiedValue(t)\bigr| \le \epsilon_{\mathrm{tol}}\,\Supply(t)$, где $\epsilon_{\mathrm{tol}} = 0.05$ поглощает задержку обнаружения.
    \item[(I4) Кредитный лимит.] $b_t(m) \ge -L\bigl(r_t(m)\bigr)$ для всех $m,t$.
    \item[(I5) Фундаментальный курс.] $P_t = \bigl(F_t + \mu\,\ExtRev_t\bigr)/\Supply(t)$ всегда вычисляется, а не назначается.
    \item[(I6) Непередаваемая репутация.] Функции переноса $\R$ не существует; $\R$ только выпускается (за заслуги) и сжигается (за нарушения).
    \item[(I7) Универсальность правил.] Любая функция зависит только от наблюдаемого состояния $(b,r,d)$, а не от личности участника.
\end{description}

Допустимые функции перехода --- \textsc{Join}, \textsc{Leave}, \textsc{Transact}, \textsc{Repay}, \textsc{Default}, \textsc{Convert}, \textsc{Invest}, а также производные периодические операции; каждая обязана сохранять (I1)--(I7).

\section{Кредит и доверие}\label{sec:confidence}

\subsection{Кредитный лимит}

\begin{equation}\label{eq:limit}
L(r) = L_0\bigl(1 + \alpha \ln(1+r)\bigr), \qquad L_0 = 100~\V,\ \ \alpha = 0.5.
\end{equation}
Логарифм ограничивает экспозицию: участник с $r = 10^4$ берёт в долг не более $\approx 560~\V$ --- лишь в $5{,}6$ раза больше лимита новичка (Таблица~\ref{tab:limit}). Поэтому ни один агент не способен угрожать платёжеспособности своим дефолтом.

\begin{table}[h]
\centering
\begin{tabular}{lrrrrrr}
\toprule
$r$ & $0$ & $1$ & $10$ & $100$ & $1000$ & $10000$\\
\midrule
$L(r)$ & $100$ & $135$ & $220$ & $330$ & $445$ & $560$\\
\bottomrule
\end{tabular}
\caption{Кредитный лимит \eqref{eq:limit} как функция репутации.}
\label{tab:limit}
\end{table}

\subsection{Доверие}

Аксиома~\ref{ax:price} запрещает механизму назначать цены; вместо этого он оценивает, насколько заявленной цене можно доверять, прежде чем выпускать под неё деньги. Для транзакции $\tau=(A,B,N,c)$ с исполнителем $A$, получателем $B$, заявленной стоимостью $N$ и категорией $c$,
\begin{equation}
\confidence(\tau) = w_a s_a(\tau) + w_r s_r(\tau) + w_x s_x(\tau), \qquad w_a+w_r+w_x=1 .
\end{equation}
\emph{Автооценка} сравнивает $N$ с недавним распределением цен в категории. При выборочном среднем $\mu_c$ и стандартном отклонении $\sigma_c$ за $90$-дневное окно и $z = (N-\mu_c)/\sigma_c$,
\begin{equation}\label{eq:autoscore}
s_a(\tau) = \begin{cases} 1, & |z| \le 1,\\[3pt] \exp\!\left(-\dfrac{(|z|-1)^2}{2}\right), & |z| > 1. \end{cases}
\end{equation}
Так $s_a = 0{,}61,\,0{,}14,\,0{,}01$ при $z = 2,3,4$. \emph{Оценка ревью} $s_r$ --- медиана голосов случайно выбранных ревьюеров (устойчива к выбросам), а \emph{оценка аудита} $s_x$ равна $1$, если стохастический аудит не выявил нарушения. Ревью запускается, когда $|z|>1$, когда $N$ превышает лимит автоодобрения $N_{\mathrm{auto}} = 50\sqrt{r_A r_B + 1}$, когда одна из сторон под наблюдением или когда транзакцию выбрал $3\%$-й случайный аудит. Эмиссия происходит только при $\confidence(\tau) \ge \theta_{\min} = 0{,}6$, согласно инварианту (I2).

\subsection{Здоровье рынка и антимонополия}

Аксиома~\ref{ax:price} помещает нечестность не в «неправильную» цену, а в отсутствие альтернатив; поэтому механизм следит за структурным здоровьем рынка, а не за отдельными ценами. Для каждой категории услуг он отслеживает индекс концентрации Херфиндаля--Хиршмана, двустороннюю концентрацию контрагентов участника и коротко-цикловую взаимность (признак отмывочных сделок). При пометке категории или кластера ответ градуирован и не касается цен: сниженные кредитные лимиты для замешанного кластера, повышенная частота аудита, приостановка компенсаторной эмиссии по подозрительным транзакциям и, в пределе, вынесение на голосование сообщества. Тем самым механизм защищает \emph{условия} честного обмена --- прозрачность и наличие альтернатив --- а не диктует условия.

\section{Эмиссия и коэффициент компенсации}\label{sec:emission}

\subsection{Механика}

В \textsc{Transact} кредитная часть равна $N_{\mathrm{credit}} = \max\bigl(0,\ N - \max(0,b_B)\bigr)$; исполнителю зачисляется $N$, а баланс $B$ уменьшается на $N$, открывая долг $N_{\mathrm{credit}}$ при $b_B < N$ (чистый перевод при $N_{\mathrm{credit}}=0$). Одновременно \emph{компенсаторная эмиссия} поступает в эскроу,
\begin{equation}\label{eq:comp}
\Delta E = \varepsilon(t)\, N_{\mathrm{credit}} .
\end{equation}
Эскроу раскрывается пропорционально при погашении и \emph{сжигается} при дефолте. Раскрытые суммы распределяются между участниками пропорционально активному обороту, где
\[
\Turnover(m,t) = \sum_{\tau:\ t-\tau<90d,\ m\in\tau} N(\tau)\,\confidence(\tau),
\]
через
\begin{equation}
\mathrm{Share}(m,t) = \frac{\Turnover(m,t) + \xi}{\sum_{m'}\Turnover(m',t) + n_t\,\xi}, \qquad \xi = 1,
\end{equation}
так что «дивиденд всем» вознаграждает активный вклад, а не праздное накопление. Коэффициент динамичен в наблюдаемой ставке дефолтов $\delta_t$:
\begin{equation}\label{eq:epsilon}
\varepsilon(t) = \max\!\Bigl(0,\ \min\bigl(\varepsilon_{\max},\ K^* - 1 - \kappa\,(\delta_t - \delta^*)\bigr)\Bigr),
\end{equation}
со значениями по умолчанию $K^* = 1{,}5$, $\delta^* = 0{,}05$, $\kappa = 2$, $\varepsilon_{\max}=0{,}95$. Тем самым $\varepsilon$ изменяется от $0{,}60$ при $\delta=0$ до $0$ при $\delta = 0{,}30$.

\subsection{Устойчивость при дефолтах}

\begin{theorem}[Ограниченность массы при дефолтах]\label{thm:stability}
Рассмотрим стационарный режим с постоянной ставкой дефолтов $\delta \le \tfrac12$ и $\varepsilon$ из \eqref{eq:epsilon}. Пусть $C_{\mathrm{tot}}$ --- суммарный подтверждённый вклад $\sum_\tau C(\tau)$, а $P_{\mathrm{net}} = (1-\delta)\,C_{\mathrm{tot}}$ --- чистый вклад, переживший дефолты. Если дефолтная эмиссия --- и её принципал, и её компенсация --- удаляется при ремедиации (операции \textnormal{\textsc{Default}} и устранения мошенничества), то, с точностью до ограниченного генезис-пула,
\[
\frac{\Supply}{P_{\mathrm{net}}} \le K^*.
\]
\end{theorem}

\begin{proof}
На единицу подтверждённого вклада механизм выпускает одну единицу принципала и при погашении раскрывает $\varepsilon$ единиц эскроу-компенсации --- итого $(1+\varepsilon)$ единиц пережившей массы; для дефолтной единицы компенсация сжигается, а принципал удаляется при ремедиации, не давая ничего. При доле дефолтов $\delta$ имеем $\Supply = (1-\delta)C_{\mathrm{tot}}(1+\varepsilon)$, тогда как $P_{\mathrm{net}} = (1-\delta)C_{\mathrm{tot}}$, так что $\Supply/P_{\mathrm{net}} = 1+\varepsilon$. По \eqref{eq:epsilon}, $\varepsilon \le K^*-1$, откуда $\Supply/P_{\mathrm{net}} \le K^*$; при $\delta > \delta^*$ слагаемое $-\kappa(\delta-\delta^*)$ делает $\varepsilon$ строго меньше, ужесточая границу.
\end{proof}

\begin{remark}
Теорема~\ref{thm:stability} --- это тормоз инфляции: с ростом фактических дефолтов $\varepsilon$ падает автоматически, а при $\delta = 0{,}30$ компенсация отключается полностью (строгая консервация). Система не предполагает полного возврата; она ограничивает отношение денег к ценности при любом $\delta \le \tfrac12$. Гипотеза о ремедиации выделяет враждебный (мошеннический) случай; добросовестный дефолт, напротив, сохраняет принципал, обеспеченный реально выполненной работой, которая затем учитывается в чистой ценности, так что граница продолжает выполняться (формализовано без какой-либо гипотезы о ремедиации в Предложении~\ref{prop:genstab}).
\end{remark}

Генезис-грант присоединяющемуся участнику, $g_0 = \min\bigl(100,\ G_t/\mathbb{E}[\text{новые участники}]\bigr)$, также самоограничен: он сжимается, если рост опережает приток в пул, так что субсидия не создаётся \emph{из ничего}.

\subsection{Разрешение долга и реактивация}

Просроченный долг влечёт градуированные санкции, а не исключение: на $30$-й, $60$-й и $90$-й день должник сталкивается с нарастающим убыванием репутации и, в конечном счёте, с заморозкой (обнуление баланса и репутации, сжигание эскроу) вместо изгнания. Замороженный участник может реактивироваться, погасив непогашенный долг и пройдя испытательный период, ограниченный мелкими транзакциями. Это регенеративное обращение вытекает из Аксиомы~\ref{ax:create}: цель --- восстановить продуктивное участие, а не наказать.

\section{Репутация и управление}

Репутация накапливается вогнутыми (активность) и линейными (качество) слагаемыми и сжигается за провалы:
\begin{align}
\Delta \R &= \beta_1 \ln(1 + n_{\mathrm{tx}}) + \beta_2\, a_{\mathrm{audit}} + \beta_3 \ln(1 + \text{стаж}/30) + \beta_4\, n_{\mathrm{disputes}},\\
\Delta \R^{-} &= \gamma_1 f_{\mathrm{review}} + \gamma_2 f_{\mathrm{dispute}} + \gamma_3\, \mathrm{DefaultPenalty} + \gamma_4\, \mathrm{Fraud},
\end{align}
где $\gamma_4 = 100$: доказанное мошенничество стирает годы репутации. Вогнутость по массовым действиям не даёт «фармить» репутацию количеством.

Голосовая мощь равна $\sqrt{\R(m)}$. Стратегические решения требуют \emph{двойного кворума}: репутационного кворума $\sum_{\text{голосующих}}\sqrt{r}\,\big/\sum_{M}\sqrt{r} > \theta$ и долевого кворума $\sum_{\text{голосующих}} \max(0,b) \,\big/ \sum_{M}\max(0,b) > \theta$, где $\theta = 0{,}5$ ($0{,}66$ для конституционных изменений). Поскольку $\R$ и $\V$ распределены по-разному, двойной кворум одновременно блокирует захват активистами и китами. Выделенный набор положений \emph{конституционен} --- семь инвариантов, непередаваемость $\R$ и запреты на premine, доли учредителей и правила, специфичные для класса --- и изменяем только референдумом сверхбольшинства, но не обычным параметрическим управлением.

\begin{proposition}[Устойчивость к капиталу и Sybil]\label{prop:sybil}
При инварианте (I6) атакующий с произвольным капиталом получает ноль репутационных голосов, поскольку $\R$ непередаваема. Sybil-атакующий, управляющий $k$ аккаунтами с базовой репутацией $r_0$, получает $k\sqrt{r_0}$ репутационных голосов; чтобы достичь цели в $q$ голосов, требуется $k = q^2/r_0$ аккаунтов. При $r_0 = 0{,}5$ для $q = 100$ голосов нужно $2\times 10^4$ аккаунтов, каждый из которых должен пройти проверку личности и вести настоящую активность, чтобы заработать хоть какую-то $\R$.
\end{proposition}

\begin{proof}
Первое утверждение немедленно следует из (I6): капитал покупает $\V$, но не $\R$, а голосовая мощь зависит только от $\R$. Для второго: аддитивность $\sqrt{\R}$ по одинаково наделённым аккаунтам даёт $k\sqrt{r_0}$; полагая $k\sqrt{r_0} = q$, получаем $k = q^2/r_0$. Подставляя $r_0=0{,}5$, $q=100$, получаем $k = 2\times 10^4$.
\end{proof}

\section{Ценообразование и фонд ликвидности}

Курс задан (I5): $P_t = (F_t + \mu\,\ExtRev_t)/\Supply(t)$ с $\mu = 12$ (мультипликатор «цена/прибыль», капитализирующий годовую внешнюю выручку).

\begin{theorem}[Нейтральность операций фонда]\label{thm:neutrality}
\textnormal{\textsc{Invest}} (внесение $\Delta U$ USDC за $\Delta V = \Delta U/P_t$ единиц $\V$) и \textnormal{\textsc{Convert}} (продажа $\Delta V$ единиц $\V$ за $\Delta U = \Delta V\,P_t$) оставляют курс неизменным: $P_{t+1} = P_t$.
\end{theorem}

\begin{proof}
Для \textsc{Invest} $F_{t+1} = F_t + \Delta U$ и $S_{t+1} = S_t + \Delta U/P_t$, где $S=\Supply$. Используя $F_t = P_t S_t - \mu\,\ExtRev_t$,
\[
P_{t+1} = \frac{F_t + \Delta U + \mu\,\ExtRev_t}{S_t + \Delta U/P_t}
= \frac{P_t S_t + \Delta U}{S_t + \Delta U/P_t}
= \frac{P_t(P_t S_t + \Delta U)}{P_t S_t + \Delta U} = P_t.
\]
Аргумент для \textsc{Convert} симметричен: $F_{t+1} = F_t - \Delta V\,P_t$ и $S_{t+1} = S_t - \Delta V$.
\end{proof}

Таким образом, курс реагирует только на фундаментальные величины ($F$ и $\ExtRev$); никакая операция фонда им не манипулирует. Кредитная эмиссия, напротив, разбавляет курс в моменте: $P_{t+1} = P_t\, S_t / \bigl(S_t + N_{\mathrm{credit}}(1+\varepsilon)\bigr) < P_t$, и восстанавливается лишь ростом $\ExtRev$. Следовательно, внутренний оборот сам по себе не создаёт внешнего богатства --- это центральный честный вывод механизма.

\paragraph{Устойчивость к банк-рану.} Пусть покрытие $\rho_t = F_t/(P_t S_t)$. При $\rho_t < \rho^* = 0{,}3$ конвертации ставятся в очередь и дисконтируются: $P_{\mathrm{actual}} = P_t \min(1, \rho_t/\rho^*)$. Дисконт и защищает фонд, и стимулирует контрциклические инвестиции на дне, восстанавливая $\rho$; он также накладывает структурный потолок эмиссии $\Delta S_{\max} = (F - F_{\min})/P_{\mathrm{target}}$.

\paragraph{Распределение прибыли.} Внешняя прибыль клуба делится каждый период: большая часть удерживается в фонде (повышая $P$ через (I5)), взвешенный по активности дивиденд выплачивается участникам, часть пополняет генезис-пул, а небольшая доля финансирует операционные расходы. Модераторы, арбитры, учредители и инвесторы --- все получают выплаты в $\V$ на одинаковых условиях: ни один класс не имеет требования старше судьбы токена или изолированного от неё. Это операционное содержание Аксиомы~\ref{ax:symmetry}: когда клуб выигрывает, держатели выигрывают вместе; когда проигрывает --- вместе проигрывают.

\section{Стационарность и условие выживания}

Курс стационарен, когда рост массы совпадает с ростом активов: $\dot S/S = \dot{\mathcal{A}}/\mathcal{A}$ с $\mathcal{A} = F + \mu\,\ExtRev$. Это разбивает динамику на режим роста, стабильный режим и режим падения.

\begin{proposition}[Условие выживания]\label{prop:survival}
Неубывающий курс требует темпа роста внешней выручки, удовлетворяющего
\[
\lambda_{\mathrm{rev}} \ \ge\ \frac{1}{\mu}\Bigl(\mathrm{EmissionRate}\cdot P - \lambda_{\mathrm{inv}}\Bigr),
\]
где $\lambda_{\mathrm{inv}}$ --- темп притока инвестиций.
\end{proposition}

\begin{proof}[Набросок]
Поскольку $P = \mathcal{A}/S$, дифференцирование даёт $\dot P = (\dot{\mathcal{A}} - P\dot S)/S$, так что неубывающий курс ($\dot P \ge 0$) эквивалентен $\dot{\mathcal{A}} \ge P\,\dot S$. Подставляя $\dot{\mathcal{A}} = \lambda_{\mathrm{inv}} + \mu\lambda_{\mathrm{rev}}$ и $\dot S = \mathrm{EmissionRate}$ (за вычетом сожжённого), получаем $\lambda_{\mathrm{inv}} + \mu\lambda_{\mathrm{rev}} \ge \mathrm{EmissionRate}\cdot P$; решая относительно $\lambda_{\mathrm{rev}}$, приходим к заявленной границе.
\end{proof}

Например, при $S = 10^5$, $P = 1$, $\mu = 12$, эмиссии $5000~\V/\text{мес}$ и инвестициях $2000~\text{USDC}/\text{мес}$ клуб должен зарабатывать снаружи не менее $250~\text{USDC}/\text{мес}$. Поэтому бесконечно падающий курс при отсутствии внешней выручки --- математическая необходимость, а не дефект дизайна.

\section{Дополнительные теоретические результаты}\label{sec:further}

Результаты ниже усиливают Теорему~\ref{thm:stability}, снимая гипотезу о ремедиации, характеризуют долгосрочный курс, показывают, что правило дисконтированного вывода не может исчерпать фонд иначе как в пределе полного выхода, и дают достаточное условие невыгодности мошенничества.

\subsection{Устойчивость без гипотезы о ремедиации}

\begin{proposition}[Обобщённая граница при смешанных дефолтах]\label{prop:genstab}
Разобьём эмиссию в стационарном режиме на погашенные кредиты (доля $1-\delta$), добросовестные дефолты (доля $\delta_g$; услуга оказана, поэтому принципал сохраняется и обеспечен ценностью, а эскроу сжигается) и обнаруженное мошенничество (доля $\delta_f$; принципал и компенсация сжигаются), где $\delta=\delta_g+\delta_f$. Считая выполненную работу ценностью, положим $P_{\mathrm{net}}=(1-\delta_f)\,C_{\mathrm{tot}}$. Тогда без каких-либо допущений о ремедиации принципала добросовестных дефолтов
\[
\frac{\Supply}{P_{\mathrm{net}}}\ \le\ 1+\varepsilon\ \le\ K^*.
\]
\end{proposition}

\begin{proof}
Погашенная единица вносит $1+\varepsilon$ в массу, добросовестный дефолт --- $1$ (принципал сохранён, эскроу сожжён), обнаруженное мошенничество --- $0$; поэтому $\Supply=C_{\mathrm{tot}}\bigl[(1-\delta)(1+\varepsilon)+\delta_g\bigr]$. Используя $\delta-\delta_f=\delta_g$,
\[
\Supply-(1+\varepsilon)P_{\mathrm{net}}
= C_{\mathrm{tot}}\bigl[(1-\delta)(1+\varepsilon)+\delta_g-(1+\varepsilon)(1-\delta_f)\bigr]
= C_{\mathrm{tot}}\bigl[\delta_g-(1+\varepsilon)\delta_g\bigr]
= -\,C_{\mathrm{tot}}\,\varepsilon\,\delta_g\ \le\ 0 .
\]
Значит $\Supply\le(1+\varepsilon)P_{\mathrm{net}}$, а $\varepsilon\le K^*-1$ из \eqref{eq:epsilon} даёт границу. При $\delta_g=0$ получаем Теорему~\ref{thm:stability}.
\end{proof}

\subsection{Долгосрочный курс}

\begin{proposition}[Долгосрочный курс]\label{prop:longrun}
Пусть агрегат активов и масса растут с постоянными темпами $g=\lambda_{\mathrm{inv}}+\mu\lambda_{\mathrm{rev}}$ и $e=\mathrm{EmissionRate}>0$ (за вычетом сожжённого), так что $\mathcal{A}(t)=\mathcal{A}_0+gt$ и $S(t)=S_0+et$. Тогда $P(t)=\mathcal{A}(t)/S(t)$ монотонен и
\[
P(t)\ \xrightarrow[t\to\infty]{}\ \frac{g}{e}=\frac{\lambda_{\mathrm{inv}}+\mu\lambda_{\mathrm{rev}}}{\mathrm{EmissionRate}} .
\]
Более того, $P$ неубывает тогда и только тогда, когда $P_0\le g/e$, что и есть условие выживания из Предложения~\ref{prop:survival}.
\end{proposition}

\begin{proof}
$P'(t)=\dfrac{g\,S(t)-e\,\mathcal{A}(t)}{S(t)^2}=\dfrac{gS_0-e\mathcal{A}_0}{S(t)^2}$, чей знак совпадает со знаком $g/e-\mathcal{A}_0/S_0=g/e-P_0$ и постоянен по $t$; значит $P$ монотонен. При $t\to\infty$ имеем $(\mathcal{A}_0+gt)/(S_0+et)\to g/e$. Наконец, $P_0\le g/e$ переписывается как $\lambda_{\mathrm{inv}}+\mu\lambda_{\mathrm{rev}}\ge\mathrm{EmissionRate}\cdot P_0$.
\end{proof}

\subsection{Фонд не исчерпывается набегом}

\begin{theorem}[Неисчерпаемость фонда]\label{thm:norun}
При набеге на вывод, целиком происходящем в дисконтированном режиме ($\rho<\rho^*$, $\ExtRev$ постоянна, без эмиссии), в непрерывном пределе фонд и масса подчиняются
\[
\frac{F}{F_0}=\Bigl(\frac{S}{S_0}\Bigr)^{1/\rho^*}.
\]
Следовательно, $F>0$ при любом $S>0$: фонд исчерпывается лишь в пределе полного выхода ($S\to0$).
\end{theorem}

\begin{proof}
Пусть $x$ --- объём сконвертированных $\V$, так что $dS/dx=-1$. В дисконтированном режиме выплата равна $P_{\mathrm{actual}}=P\,\rho/\rho^*$; используя $P=\mathcal{A}/S$ и $\rho=F/\mathcal{A}$, получаем $P_{\mathrm{actual}}=F/(\rho^*S)$, откуда $dF/dx=-F/(\rho^*S)$. Значит $dF/dS=F/(\rho^*S)$ --- разделяющееся уравнение с решением $\ln(F/F_0)=\rho^{*-1}\ln(S/S_0)$, то есть $F=F_0(S/S_0)^{1/\rho^*}$. Поскольку $1/\rho^*>0$, $F$ и $S$ обращаются в ноль вместе.
\end{proof}

\begin{remark}
Поскольку выплата $P_{\mathrm{actual}}=F/(\rho^*S)$ убывает вместе с $F$, предельная выручка держателя стремится к нулю по мере углубления набега, снимая стимул продолжать конвертацию; вместе с теоремой это структурный ответ механизма на панику.
\end{remark}

\subsection{Невыгодность мошенничества (достаточное условие)}

\begin{proposition}[Невыгодность мошенничества]\label{prop:fraud}
Рассмотрим мошеннический кластер, который получает $W>0$, если успевает конвертировать извлечённую ценность в USDC до обнаружения, но теряет залог $L>0$ (сожжённые балансы и репутация), если обнаружен раньше. Пусть обнаружение происходит независимо каждый период с вероятностью не менее $p$, а новые аккаунты не могут конвертировать в течение lock-up из $\tau_{\mathrm{lock}}$ периодов. Тогда ожидаемая прибыль удовлетворяет
\[
\mathbb{E}[\pi]\ \le\ (1-p)^{\tau_{\mathrm{lock}}}\,W-\bigl(1-(1-p)^{\tau_{\mathrm{lock}}}\bigr)L,
\]
что неположительно, как только $(1-p)^{\tau_{\mathrm{lock}}}\le L/(W+L)$.
\end{proposition}

\begin{proof}
До истечения $\tau_{\mathrm{lock}}$ конвертация невозможна, поэтому извлечь $W$ можно, лишь оставаясь необнаруженным не менее $\tau_{\mathrm{lock}}$ периодов --- событие вероятности не более $(1-p)^{\tau_{\mathrm{lock}}}$. Ограничивая выигрыш величиной $W$ на этом событии и потерю величиной $L$ на дополнении, получаем неравенство; требуя $\mathbb{E}[\pi]\le0$ и решая, приходим к $(1-p)^{\tau_{\mathrm{lock}}}\le L/(W+L)$.
\end{proof}

\begin{remark}
При lock-up по умолчанию ($\tau_{\mathrm{lock}}=60$ дней) и любой ненулевой периодической интенсивности обнаружения от $3\%$-го стохастического аудита и мониторинга концентрации из раздела~\ref{sec:confidence} порог выполняется при реалистичных $W/L$: lock-up вынуждает атакующего держать необеспеченную, обнаружимую позицию достаточно долго, чтобы обнаружение доминировало. Полный теоретико-игровой оптимум по адаптивным стратегиям остаётся открытым (раздел~\ref{sec:limitations}).
\end{remark}

\section{Эмпирическая валидация}

\subsection{Метод}

Мы реализовали полный механизм как детерминированный агентный симулятор с инъектируемым сеяным источником псевдослучайности, так что один seed воспроизводит прогон побитово. Устойчивость оценивается пятью критериями (V2): (C1) курс не достигает нуля; (C2) покрытие фонда $\rho \ge 0{,}05$ хотя бы на $95\%$ тиков; (C3) заморожено не более $20\%$ участников; (C4) неотрицательная масса; (C5) инварианты чисты. Мы сообщаем (C1)--(C4) как \emph{экономически устойчивые} отдельно от (C5), поскольку проверка (I3) имеет известный дрейф на жизненном цикле авто-погашения/распределения эскроу. Протокол включает: воспроизведение эталонных сценариев (V1); Монте-Карло по $4$ сценариям $\times\ 30$ seed-ов (V3); развёртку из $192$ точек по $(K^*,\mu,\delta^*,\rho^*)$ (V4); стресс-тесты с комбинированными атаками (V5); и анализ режимов (V6).

\subsection{Результаты}

На $315$ прогонах ни один не привёл к экономическому коллапсу (Таблица~\ref{tab:results}). Развёртка по параметрам была устойчива во всех $192$ точках, что указывает на неломкость механизма к выбору параметров. Комбинированные стресс-тесты «мошенничество плюс банк-ран» и «одновременный выход $70\%$» оставались экономически устойчивыми за счёт очереди вывода и сжигания эскроу, а три предсказанных режима курса были воспроизведены.

\begin{table}[h]
\centering
\begin{tabular}{lrr}
\toprule
Протокол & Прогонов & Коллапсов\\
\midrule
Монте-Карло (V3) & $120$ & $0$\\
Развёртка параметров (V4) & $192$ & $0$\\
Стресс-тесты (V5) & $3$ & $0$\\
\midrule
Всего & $315$ & $0$\\
\bottomrule
\end{tabular}
\caption{Сводка валидации. Полные сырые результаты сопровождают выложенный код.}
\label{tab:results}
\end{table}

\section{Ограничения}\label{sec:limitations}

Эмпирическая устойчивость необходима, но не достаточна. Наши прогоны покрывают $1$--$3$ месяца виртуального времени; медленные эффекты на многолетних горизонтах не проверены. Агенты статистические, а не стратегические, поэтому адаптация, сговор и политическая самоорганизация реальных участников вне охвата, как и теоретико-игрово оптимальный противник (мы тестируем лишь фиксированные стратегии атак). Правовые и налоговые вопросы --- юрисдикция, идентификация/KYC и налогообложение микротранзакций --- лежат вне модели. Поэтому механизм представляет собой структурно надёжный фундамент, поведение которого с реальными людьми ещё предстоит установить.

\section{Заключение}

Kredo показывает, что связная этическая позиция --- вознаграждение за создание, симметрия судьбы, прозрачность и власть от вклада --- может быть скомпилирована в точный механизм, чьи заложенные свойства доказуемы и машинно-проверяемы. Коэффициент компенсаторной эмиссии ограничивает инфляцию при дефолтах, операции фонда нейтральны к курсу, управление устойчиво к захвату капиталом и Sybil, а минимальное условие по внешней выручке управляет выживанием. Детерминированный симулятор подтверждает структурную устойчивость на $315$ прогонах. Открытая публикация приглашает к воспроизведению, состязательному анализу и расширению в сторону небольшого живого запуска.

\section*{Доступность данных и кода}

Спецификация механизма, полные выводы, симулятор, протоколы экспериментов и сырые результаты открыто доступны по адресу \url{https://github.com/fedorello/kredo-model} под лицензией MIT.

\begin{thebibliography}{99}
\bibitem{gesell1916} S.~Gesell. \emph{The Natural Economic Order}. Self-published, 1916 (English translation, Peter Owen, 1958).
\bibitem{greco2009} T.~H.~Greco. \emph{The End of Money and the Future of Civilization}. Chelsea Green, 2009.
\bibitem{stodder2009} J.~Stodder. Complementary credit networks and macroeconomic stability: Switzerland's Wirtschaftsring. \emph{Journal of Economic Behavior \& Organization}, 72(1):79--95, 2009.
\bibitem{circles2020} Circles~UBI. \emph{Circles: A Basic Income Money System} (whitepaper). 2020.
\bibitem{gooddollar2021} GoodDollar. \emph{GoodDollar: A Distributed Basic Income} (whitepaper), v2. 2021.
\bibitem{nakamoto2008} S.~Nakamoto. \emph{Bitcoin: A Peer-to-Peer Electronic Cash System}. 2008.
\bibitem{lalley2018} S.~P.~Lalley and E.~G.~Weyl. Quadratic voting: how mechanism design can radicalize democracy. \emph{AEA Papers and Proceedings}, 108:33--37, 2018.
\bibitem{buterin2019} V.~Buterin, Z.~Hitzig, and E.~G.~Weyl. A flexible design for funding public goods. \emph{Management Science}, 65(11), 2019.
\bibitem{douceur2002} J.~R.~Douceur. The Sybil attack. In \emph{Proc.\ IPTPS}, LNCS 2429, pp.~251--260, 2002.
\bibitem{bonabeau2002} E.~Bonabeau. Agent-based modeling: methods and techniques for simulating human systems. \emph{PNAS}, 99(suppl.~3):7280--7287, 2002.
\end{thebibliography}

\end{document}
